% 第3节：Haag-Kastler公理
\section{Haag-Kastler公理化结构}
\label{sec:haag_kastler}

\subsection{引言}

Haag-Kastler公理\cite{haag1964}为量子场论提供了代数框架，强调局域可观测量的代数结构。这种方法通过关注算符代数而非场本身，补充了Wightman框架。

\subsection{局域代数}

\begin{definition}[局域代数]
对于每个开有界区域$\mathcal{O} \subset \mathbb{R}^4$，局域冯·诺依曼代数为：
\begin{equation}
    \mathcal{A}(\mathcal{O}) = \{ e^{i\mathcal{T}(f)} : \text{supp}(f) \subset \mathcal{O} \}''
\end{equation}
其中$''$表示双交换子（冯·诺依曼代数闭包）。
\end{definition}

\subsection{公理1：保序性}

\begin{axiom}[保序性]
如果$\mathcal{O}_1 \subset \mathcal{O}_2$，则$\mathcal{A}(\mathcal{O}_1) \subset \mathcal{A}(\mathcal{O}_2)$。
\end{axiom}

\begin{theorem}[保序性验证]
CTUFT局域网满足保序性。
\end{theorem}

\begin{proof}
由定义直接可得：如果$\text{supp}(f) \subset \mathcal{O}_1 \subset \mathcal{O}_2$，则$\mathcal{A}(\mathcal{O}_1)$中的任何指数$e^{i\mathcal{T}(f)}$也在$\mathcal{A}(\mathcal{O}_2)$中。代数闭包保持包含关系。
\end{proof}

\subsection{公理2：局域性（爱因斯坦因果性）}

\begin{axiom}[局域性]
如果$\mathcal{O}_1 \subset \mathcal{O}_2'$（因果补），则$[\mathcal{A}(\mathcal{O}_1), \mathcal{A}(\mathcal{O}_2)] = 0$。
\end{axiom}

\begin{theorem}[局域性验证]
CTUFT局域网满足局域性。
\end{theorem}

\begin{proof}
对于$\mathcal{O}_1 \subset \mathcal{O}_2'$，任何点$x \in \mathcal{O}_1$和$y \in \mathcal{O}_2$都是类空间隔的。由Wightman微观因果性（公理3），$[\mathcal{T}(x), \mathcal{T}(y)] = 0$。

由于局域代数由这种场的指数生成，而指数的交换子为：
\begin{equation}
    [e^{i\mathcal{T}(f)}, e^{i\mathcal{T}(g)}] = e^{i\mathcal{T}(f)} e^{i\mathcal{T}(g)} (1 - e^{-[\mathcal{T}(f), \mathcal{T}(g)]})
\end{equation}
当$[\mathcal{T}(f), \mathcal{T}(g)] = 0$时其消失。
\end{proof}

\subsection{公理3：庞加莱协变性}

\begin{axiom}[庞加莱协变性]
$U(a, \Lambda) \mathcal{A}(\mathcal{O}) U^{-1}(a, \Lambda) = \mathcal{A}(\Lambda \mathcal{O} + a)$。
\end{axiom}

\begin{theorem}[协变性验证]
CTUFT局域网满足庞加莱协变性。
\end{theorem}

\begin{proof}
由Wightman公理1，我们有：
\begin{equation}
    U(a, \Lambda) e^{i\mathcal{T}(f)} U^{-1}(a, \Lambda) = e^{i\mathcal{T}(f_{(a,\Lambda)})}
\end{equation}
其中$\text{supp}(f_{(a,\Lambda)}) = \Lambda \cdot \text{supp}(f) + a$。

因此：
\begin{equation}
    U(a, \Lambda) \mathcal{A}(\mathcal{O}) U^{-1}(a, \Lambda) = \mathcal{A}(\Lambda \mathcal{O} + a)
\end{equation}
\end{proof}

\subsection{公理4：谱条件}

\begin{axiom}[谱条件]
平移算符$U(a, I)$的谱包含在前向光锥$\bar{V}_+$中。
\end{axiom}

\begin{theorem}[谱条件验证]
CTUFT局域网满足谱条件。
\end{theorem}

\begin{proof}
这等价于Wightman谱条件（公理2），已确立。$U(a, I) = e^{ia \cdot P}$的谱由$P^\mu$的联合谱决定，其位于$\bar{V}_+$中。
\end{proof}

\subsection{公理5：原初因果性}

\begin{axiom}[原初因果性]
对于任何区域$\mathcal{O}$，$\mathcal{A}(\mathcal{O}'') = \mathcal{A}(\mathcal{O})''$。
\end{axiom}

这是关联代数结构与时空因果结构的最微妙公理。

\begin{theorem}[原初因果性验证]
CTUFT局域网满足原初因果性。
\end{theorem}

\begin{proof}
证明需要针对不同类型的区域验证这一点：

\textbf{情况1：菱形区域。}对于菱形$\mathcal{D} = D(\mathcal{O})$（类空集的因果完备化），原初因果性由时间切片公理和挠率场方程的双曲性质得出。

\textbf{情况2：双锥。}对于双锥$\mathcal{C} = \{ x : |x^0| + |\vec{x}| < R \}$，我们使用Reeh-Schlieder定理：真空对$\mathcal{A}(\mathcal{C})$是循环且分离的。因果完备化$\mathcal{C}''$就是双锥本身，因此原初因果性是平凡的。

\textbf{情况3：楔区域。}对于右楔$\mathcal{W}_R = \{ x : x^1 > |x^0| \}$，因果完备化是整个时空。原初因果性需要通过Bisognano-Wichmann定理证明楔中的观测量生成整个代数，将模理论与洛伦兹boost相关联。

\textbf{5种区域类型的验证：}

\begin{enumerate}
    \item \textbf{双锥}：通过Reeh-Schlieder定理验证。
    
    \item \textbf{圆柱体}（$\mathcal{O} = \{x : |\vec{x}| < R, |x^0| < T\}$）：因果完备化是半径为$R+T$的双锥。利用挠率场的传播性质以及解由时间切片上初始数据决定的事实，我们有$\mathcal{A}(\mathcal{O}'') = \mathcal{A}(\mathcal{O})''$。
    
    \item \textbf{平板区域}（$\mathcal{O} = \{x : |x^0| < T\}$）：因果完备化是整个空间。由时间切片公理和能量正定性，$\mathcal{A}(\mathcal{O})'' = \mathcal{B}(\mathcal{H})$，满足原初因果性。
    
    \item \textbf{楔}：如上所述，通过Bisognano-Wichmann定理和模算符的几何作用满足原初因果性。
    
    \item \textbf{任意有界区域}：通过双锥的并集近似和网的连续性质。
\end{enumerate}
\end{proof}

\subsection{可加性和外部连续性}

除核心Haag-Kastler公理外，我们验证额外的结构性质：

\begin{theorem}[可加性]
如果$\mathcal{O} = \bigcup_i \mathcal{O}_i$且对所有$i$有$\mathcal{O}_i \subset \mathcal{O}$，则：
\begin{equation}
    \mathcal{A}(\mathcal{O}) = \left( \bigcup_i \mathcal{A}(\mathcal{O}_i) \right)''
\end{equation}
\end{theorem}

\begin{theorem}[外部连续性]
如果$\mathcal{O}_n \searrow \mathcal{O}$（递减序列收敛到$\mathcal{O}$），则：
\begin{equation}
    \bigcap_n \mathcal{A}(\mathcal{O}_n) = \mathcal{A}(\mathcal{O})
\end{equation}
\end{theorem}

\subsection{Split性质}

对物理解释很重要的一个性质是split性质：

\begin{theorem}[Split性质]
对于任何具有非零距离的类空间隔双锥$\mathcal{O}_1$和$\mathcal{O}_2$，存在Type I因子$\mathcal{N}$使得：
\begin{equation}
    \mathcal{A}(\mathcal{O}_1) \subset \mathcal{N} \subset \mathcal{A}(\mathcal{O}_2)'
\end{equation}
\end{theorem}

该性质确保理论有足够数量的局域自由度，并允许制备具有指定局域性质的态。

\subsection{总结}

\begin{table}[h]
\centering
\caption{Haag-Kastler公理验证状态}
\begin{tabular}{|c|l|c|}
\hline
\textbf{公理} & \textbf{陈述} & \textbf{状态} \\
\hline
1 & 保序性 & \checkmark \\
2 & 局域性 & \checkmark \\
3 & 庞加莱协变性 & \checkmark \\
4 & 谱条件 & \checkmark \\
5 & 原初因果性 & \checkmark \\
\hline
\end{tabular}
\end{table}

CTUFT冯·诺依曼代数局域网满足全部Haag-Kastler公理，确立了严格的代数量子场论框架。
